\section{Frontier RL Flywheel y límites de capacidad}

Con un entorno real, el sistema puede generar tareas, intentar acciones, verificar resultados y luego usar las trayectorias exitosas para actualizar el modelo. El ponente llama a este proceso recursive self-improvement, pero el término se malinterpreta con facilidad, como si el modelo se separara por completo de las personas y se mejorara de manera autónoma e ilimitada. Una explicación más precisa del sistema es esta: el modelo ayuda a ampliar la experimentación y la producción de datos, mientras que las personas y la infraestructura siguen encargándose de elegir los entornos, definir la verificación, controlar el despliegue, atender las anomalías y decidir qué mejoras merecen pasar a la siguiente ronda.

\begin{figure}[H]
\centering
\includegraphics[width=0.94\textwidth]{images/00-35-15-frontier-rl-flywheel.jpg}
\caption{Frontier RL flywheel: parte de hipótesis y entornos, pasa por el intento, la verificación y el aprendizaje, y genera un sistema más fuerte para la siguiente ronda.}
\figsource{Video oficial de Stanford Online, 00:35:15, `images/00-35-15-frontier-rl-flywheel.jpg`.}
\end{figure}

\begin{knowledgebox}{Lectura de la figura: la mejora recursiva necesita anclas externas}
A la izquierda está la generación de tareas, datos o hipótesis científicas; al centro, el modelo actúa en el entorno y obtiene resultados verificables; a la derecha, la experiencia útil se incorpora de vuelta a la política, las herramientas o la base de conocimiento. Las anclas externas del ciclo incluyen pruebas unitarias, experimentos físicos, aceptación por parte del usuario, reglas de seguridad y límites de costo. Sin anclas confiables, el sistema podría limitarse a optimizar métricas sustitutas y caer en un autorreforzamiento que parece rápido, pero que en realidad se aleja del objetivo.
\end{knowledgebox}

\subsection{Límite uno: cuánto más puede escalar en la práctica}

El límite empírico se ocupa de curvas observables: al aumentar los rollouts, el número de entornos, la dificultad de las tareas y la capacidad de cómputo para entrenamiento, ¿sigue subiendo la tasa de éxito?, ¿cuándo empiezan a disminuir los rendimientos marginales?, ¿aparece colapso de modos o reward hacking durante el entrenamiento? Esta pregunta puede responderse poco a poco mediante experimentos. La respuesta también varía según el dominio: las pruebas de código, las respuestas matemáticas, los experimentos de materiales y la escritura abierta tienen costos de verificación completamente distintos.

\begin{figure}[H]
\centering
\includegraphics[width=0.9\textwidth]{images/00-37-15-limits-to-rl.jpg}
\caption{El ponente plantea dos tipos de límites de RL: la meseta empírica y el límite filosófico relativo a los objetivos y los valores humanos.}
\figsource{Video oficial de Stanford Online, 00:37:15, `images/00-37-15-limits-to-rl.jpg`.}
\end{figure}

\begin{warningbox}{Lectura de la figura: los dos tipos de pregunta no se responden con la misma curva}
«¿Sigue subiendo la exactitud con más capacidad de cómputo?» es una pregunta empírica, para la que pueden diseñarse experimentos controlados; «¿qué se debe optimizar?, ¿quién tiene derecho a definir el éxito?, ¿qué capacidades no deben desplegarse?» son preguntas normativas y de gobernanza. Estas últimas no desaparecen por sí solas porque el benchmark siga subiendo. Mezclar ambas hace que la curva técnica sustituya a la elección de valores, y también que el debate sobre valores ignore la evidencia de ingeniería real.
\end{warningbox}

\subsection{Límite dos: de quién es el objetivo que representa la reward}

El límite filosófico termina por convertirse en una interfaz del sistema. Quién elige las tareas de entrenamiento, quién escribe el verifier, quién puede modificar la reward, quién asume los fallos y si el usuario sabe qué está optimizando el sistema: ninguna de estas es una cuestión accesoria que se atienda después del entrenamiento. La ventaja de RL es que incorpora la retroalimentación al aprendizaje; su riesgo proviene de ese mismo punto: si el criterio de retroalimentación es demasiado estrecho, manipulable o poco representativo, la optimización amplificará el sesgo hasta convertirlo en una conducta estable.

\begin{table}[H]
\centering
\begin{tabular}{L{0.2\textwidth}L{0.31\textwidth}L{0.39\textwidth}}
\toprule
Tipo de límite & Señal observable & Forma de respuesta \\
\midrule
Límite de cómputo & Rendimiento de rollouts, duración del entrenamiento, energía y costo & Optimización del sistema, eficiencia de muestreo, presupuesto de recursos \\
Límite de entorno & Cobertura insuficiente de tareas, desfase entre simulación y realidad & Ampliar los entornos reales, verificación por capas, auditoría humana \\
Límite de verificación & Reward hacking, filtración de pruebas, distorsión de métricas sustitutas & Varios verifiers, pruebas ocultas, repositorio de casos de fallo \\
Límite organizacional & Responsabilidad difusa, presión de lanzamiento que se impone a las señales de seguridad & Autoridad de decisión clara, análisis posterior de incidentes, condiciones de detención \\
Límite de valores & Conflicto entre grupos sobre la definición de éxito & Diseño participativo, restricciones institucionales, mecanismos de apelación \\
\bottomrule
\end{tabular}
\caption{Descomposición de «los límites de RL» en problemas distintos que pueden medirse y gobernarse.}
\end{table}

\begin{knowledgebox}{Nota de clase: primero, recorrer el ciclo por cuenta propia}
El ponente anima a los estudiantes a entender RL mediante el problem set, en lugar de juzgarlo solo a partir de narrativas generales. Definir uno mismo el entorno, escribir la reward y observar cómo la política aprovecha los huecos pone en evidencia muy pronto la distancia entre «un objetivo que parece claro» y «un objetivo que puede calcularse de forma confiable». Esta práctica también es preparedness: saber cuándo conviene confiar en la curva y cuándo hay que revisar los datos y el verifier.
\end{knowledgebox}

\subsection{Resumen de la sección}

El Frontier RL flywheel puede ampliar la producción de experiencia verificable, pero cada ronda depende de los entornos, el verifier, las decisiones humanas y la infraestructura. El límite empírico puede explorarse con curvas y experimentos; los límites filosóficos y de gobernanza requieren responsabilidades claras y deliberación pública. A continuación, la perspectiva pasa del ciclo de entrenamiento individual al mercado en su conjunto, para observar si el gasto de capital, la adopción de software y el precio del cómputo respaldan la narrativa de la «gran transformación».
